\documentclass[11pt,UTF8,fontset=fandol]{ctexart}
\usepackage{amsmath,amssymb,amsfonts}
\usepackage{graphicx,float,caption,subcaption}
\usepackage{booktabs,array,multirow,makecell}
\usepackage{geometry,enumitem,xcolor}
\usepackage[hidelinks]{hyperref}
\usepackage{url}
\geometry{margin=2.3cm}
\graphicspath{{/Users/yye/project/论文/assets/stepcoder/}}
\captionsetup{font=small,labelfont=bf}
\setlength{\parskip}{0.4em}

% 小白知识卡 / 小贴士 盒子
\newenvironment{infobox}[1]{%
  \par\medskip\noindent
  \begin{tabular}{@{}p{0.97\linewidth}@{}}
  \hline
  \rule{0pt}{2.4ex}\textbf{#1}\\
}{%
  \\\hline
  \end{tabular}\par\medskip
}

\title{\textbf{StepCoder：用编译器反馈的强化学习改进代码生成（中文精读导读）}}
\author{原作者：Shihan Dou, Yan Liu, Haoxiang Jia, Limao Xiong, Enyu Zhou, Wei Shen 等\\
Fudan NLP Lab, Fudan University（复旦大学）等机构}
\date{}

\begin{document}
\maketitle

\begin{center}
\fbox{\begin{minipage}{0.95\linewidth}
\small
\textbf{英文原标题：} StepCoder: Improve Code Generation with Reinforcement Learning from Compiler Feedback\\[2pt]
\textbf{发表：} ACL 2024（arXiv:2402.01391）\quad\textbf{引用数：} 较高（本次未取到精确数值）\\[2pt]
\textbf{代码与数据集：} \url{https://github.com/Ablustrund/APPS_Plus}\\[4pt]
\textbf{一句话定位：} 让大模型写代码时，\emph{把"代码能不能跑通单元测试"当作奖励信号来做强化学习}；并用两招——\textbf{把"从零写整段难代码"拆成"由易到难的代码补全课程"（CCCS）}、\textbf{只对真正被测试执行到的代码算梯度（FGO）}——解决"探索难"和"优化不准"两个老大难问题。
\end{minipage}}
\end{center}

\noindent\textit{说明：本文是面向 AI 应用开发实习生的\textbf{中文精读导读}（转化性学习材料），以讲解、转述、关键句翻译和公式拆解为主，并非原文逐字全译。术语首次出现时括注英文原文，模型/数据集/方法专名保留英文。}

% ======================================================================
\section{小白知识卡：读这篇前要懂的几个概念}
% ======================================================================

这篇论文把"代码生成"问题套进了"强化学习（reinforcement learning, RL）"的框架。如果你对 RL 还不太熟，先把下面几个概念过一遍，后面就好读了。

\begin{infobox}{代码生成（code generation）}
给模型一段自然语言的需求描述（比如"写一个统计子序列个数的函数"），让它自动生成能正确运行的源代码。本文用的"对错"标准很硬核：生成的代码要能通过预先准备好的\textbf{单元测试（unit test）}——也就是给定输入，输出必须和标准答案一致。
\end{infobox}

\begin{infobox}{强化学习（RL）三件套：状态 / 动作 / 奖励}
RL 是"在环境里不断试，根据奖励调整策略"。在写代码这个场景里：\textbf{动作（action）}=模型每生成一个 token（词元）；\textbf{状态（state）}=已经生成的前面那串 token；\textbf{奖励（reward）}=整段代码写完、丢进编译器跑完单元测试后，告诉你"过了 / 没过 / 报错"。负责生成的模型叫\textbf{策略模型（policy model）}。
\end{infobox}

\begin{infobox}{PPO（近端策略优化，Proximal Policy Optimization）}
目前最常用的 RL 训练算法之一，也是 RLHF（基于人类反馈的强化学习）里的主力。它让策略模型"小步快走"地朝高奖励方向更新，同时用一个 KL 惩罚项拴住它、不让它跑得离原始模型太远（防止训崩）。本文就是用 PPO 来训练代码模型的。
\end{infobox}

\begin{infobox}{稀疏奖励（sparse reward）与探索（exploration）难题}
"稀疏奖励"是指：只有把\textbf{整段}代码生成完、跑完测试，才会拿到一次奖励信号；中间每写一个 token 都没有任何反馈。需求越复杂，要写的代码越长，模型靠"瞎试"凑出一份能通过测试的代码的概率就越低——这就是"探索难"。本文的 CCCS 就是专门治这个的。
\end{infobox}

\begin{infobox}{优势函数（advantage）与信用分配（credit assignment）}
RL 里用\textbf{优势函数 $A$} 衡量"某个动作比平均水平好多少"。难点在于：一段代码里有很多 token，但单元测试可能根本没执行到其中一部分。如果对所有 token 一视同仁地算优势、算梯度，就会把"功劳/锅"算到没参与的代码头上——这叫信用分配不准。本文的 FGO 就是专门治这个的。
\end{infobox}

\begin{infobox}{课程学习（curriculum learning）}
像教小孩一样，\textbf{先学简单的、再学难的}。本文把"从零写一整段复杂代码"这件难事，先变成"我给你大半段标准答案、你只补最后一点点"的简单填空，再逐步减少给的提示、加大要补的比例，直到模型能完全独立写出整段——这就是 CCCS 的核心思路。
\end{infobox}

% ======================================================================
\section{摘要全译}
% ======================================================================

大语言模型（large language models, LLMs）的发展极大推动了代码生成领域。已有工作把强化学习（RL）与编译器反馈结合起来，借此探索 LLM 的输出空间、提升代码生成质量。然而，LLM 面对复杂人类需求时生成的代码往往很长，这使得 RL 的探索成为一大挑战；同时，由于单元测试可能覆盖不到全部复杂代码，用这些"未被执行的代码片段"来优化 LLM 是无效的。为应对这些挑战，我们提出 StepCoder——一个用于代码生成的全新 RL 框架，由两个主要组件构成：CCCS 通过把"长序列的代码生成任务"分解为"由一系列代码补全子任务组成的课程（Curriculum of Code Completion Subtasks）"来解决探索难题；FGO 则通过屏蔽（mask）未被执行的代码段，只用真正执行到的代码来优化模型，从而实现细粒度优化（Fine-Grained Optimization）。此外，我们还构建了用于 RL 训练的 APPS+ 数据集，并经过人工核验以确保单元测试的正确性。实验结果表明，我们的方法提升了探索输出空间的能力，并在相应基准上超越了当前最先进的方法。我们的数据集 APPS+ 与 StepCoder 已在线公开。

% ======================================================================
\section{逐节精读}
% ======================================================================

\subsection{引言（Introduction）：要解决什么问题}

代码生成（也叫程序合成 program synthesis）的目标，是根据自然语言描述自动生成符合要求的源代码。基于 LLM 的方法在这件事上已经表现亮眼，但在"对齐复杂人类需求"上仍有差距。

一个很有潜力的方向是\textbf{从编译器反馈中学习}：把代码丢进编译器、跑单元测试，结果直接告诉你这段程序功能上对不对。前人（如 PPOCoder、RLTF、CodeRL）已经把"单元测试通过与否"当作奖励，用 RL 来引导模型探索输出空间。但作者指出，用 RL 训代码模型有三个拦路虎：

\begin{enumerate}[leftmargin=2em]
\item \textbf{探索难}：需求越复杂，代码越长，稀疏奖励下模型很难"探"到一份能通过测试的解。
\item \textbf{优化不准}：单测覆盖不到的代码片段和奖励无关，如果拿整段代码去算梯度，优势估计就不准。
\item \textbf{数据集质量差}：现有数据集（如 APPS）存在缺输入/输出、标准答案跑不通、输出对不上等问题，影响 RL 学习。
\end{enumerate}

针对这三点，作者的贡献是：(1) 提出 StepCoder，含 CCCS 和 FGO 两个组件；(2) 构建高质量数据集 APPS+；(3) 实验证明 StepCoder 能提升探索效率与效果，超越其他方法。

\textit{关键句直译：} "CCCS makes exploration easier by breaking down the complicated goals into sub-objectives curriculum. FGO provides fine-grained optimization by only utilizing the executed code in unit tests."（CCCS 通过把复杂目标拆成"子目标课程"让探索更容易；FGO 通过只利用单元测试中被执行到的代码来提供细粒度优化。）

\subsection{动机（Motivation）：两个难题的数学直觉}

作者用 APPS 数据集里一个例子（见原文图 1）把两个难题讲清楚。他们先用\textbf{抽象语法树（abstract syntax tree, AST）}分析代码，找出其中的\textbf{条件语句（conditional statement, CS）}——也就是 \texttt{if/for} 这类会产生新分支、增加程序复杂度的语句。

\paragraph{难题一：探索问题。} 设 $P_\theta(\mathrm{CS}_i)$ 表示参数为 $\theta$ 的策略模型\textbf{正确写出第 $i$ 个条件语句}的概率。那么模型按需求正确生成整段代码的概率可以写成：
\begin{equation*}
P \propto P_o \prod_{i=1}^{3} P_\theta(\mathrm{CS}_i) \tag{1}
\end{equation*}

\noindent\textbf{逐符号拆解：}
\begin{itemize}[leftmargin=2em]
\item $P$：整段代码被正确生成的总概率（正比于右边）。
\item $P_o$：图中标注的条件语句之外、其它代码片段被写对的概率。
\item $\prod_{i=1}^{3} P_\theta(\mathrm{CS}_i)$：把每个条件语句写对的概率\textbf{连乘}起来（这个例子里有 3 个）。
\end{itemize}

直觉是：每个 $P_\theta(\mathrm{CS}_i)$ 都小于 1，\textbf{连乘}会让总概率迅速变小。需求越复杂、条件语句越多，连乘项越多，$P$ 就会\textbf{指数级地塌下去}——模型几乎不可能靠探索碰巧写对整段，这就是稀疏奖励 + 长序列下探索极难的根源。作者也指出，常见的"reward shaping（奖励塑形，更频繁地人为给奖励）"在这里行不通，因为代码必须\textbf{完整生成完}才能跑、才有奖励。

\paragraph{难题二：优化问题。} 先看标准 RL 微调的目标：把每个 token 的预测看成一个动作 $a$，更新策略模型 $\pi_\theta$ 的梯度策略目标为：
\begin{equation*}
\max_\theta\ \mathbb{E}_{(x,y)\sim D_{\pi_\theta}}\Big[\sum_t A^\pi_t \log \pi(y_t \mid y_{1:t-1}, x; \theta)\Big] \tag{2}
\end{equation*}

\noindent\textbf{逐符号拆解：}
\begin{itemize}[leftmargin=2em]
\item $x$：人类需求（任务描述）；$y$：生成的代码（一个 episode，整段轨迹）。
\item $\log \pi(y_t \mid y_{1:t-1}, x; \theta)$：在已生成前 $t-1$ 个 token 和需求 $x$ 的条件下，生成第 $t$ 个 token $y_t$ 的对数概率。
\item $A^\pi_t$：第 $t$ 个动作（token）的\textbf{优势}，由 GAE（广义优势估计，Generalized Advantage Estimator）从奖励 $r$ 算出，用来降低方差。
\item 整个式子的意思：\textbf{对优势为正的 token 提高其概率、对优势为负的降低其概率}，从而让模型更倾向于高回报的写法。
\end{itemize}

问题来了：奖励只跟\textbf{被执行到的代码}有关。原文举例：某个输入下只有 75\% 的代码被执行，那剩下 25\% 没执行的 token 跟奖励根本无关。可是式 (2) 把\textbf{所有} token 都拉进梯度计算，就会给这些无关 token 也分配优势——导致优势估计不准、优化不精确。

\subsection{方法（Method）：StepCoder 怎么做}

StepCoder 的整体框架见图~\ref{fig:overview}。先约定记号（Preliminaries）：训练集 $D = \{(x_i, y_i, u_i, e_i)\}_{i=0}^{N}$，其中 $x$ 是人类需求、$y$ 是标准答案（canonical solution）、$u$ 是单元测试样本、$e$ 是通过分析 AST 得到的\textbf{条件语句列表}（每个元素记录该语句的起始位置 $st$ 和结束位置 $en$，按 $st$ 升序排列）。标准答案 $y$ 可表示为 token 序列 $\{a_t\}_{t=0}^{T}$。

\paragraph{组件一：CCCS（代码补全子任务课程）。} 核心思路是\textbf{把"长序列探索"拆成"一串由易到难的短任务"}。具体做法：

\begin{itemize}[leftmargin=2em]
\item \textbf{从靠近终点的地方起步。} 训练早期，给模型需求 $x$ \emph{加上}标准答案的前面一大段 $x_p = \{a_i\}_{i=0}^{s^*}$ 当提示，让模型只补完剩下一小段。这里 $s^*$ 是探索的起始点。把拼好的输入记为 $x' = (x, x_p)$，模型输出轨迹 $\tau$，最终序列 $\hat{y} = (x_p, \tau)$。因为只需补一点点，探索一下就成功了，奖励信号容易拿到。
\item \textbf{奖励函数（沿用前人设置）：}
\begin{equation*}
r(x', \hat{y}) =
\begin{cases}
+1, & \text{若 } \hat{y} \text{ 通过全部单元测试}\\
-0.3, & \text{若 } \hat{y} \text{ 有任一单元测试失败}\\
-0.6, & \text{若 } \hat{y} \text{ 发生运行时错误（runtime error）}\\
-1, & \text{若 } \hat{y} \text{ 发生编译错误（compile error）}
\end{cases}
\tag{3}
\end{equation*}
直觉：通过给正分；越"早死"（编译就过不了）罚得越重，引导模型至少先写出能跑的代码。
\item \textbf{用 PPO 优化，并屏蔽提示部分。} 在优化阶段，当作提示喂进去的标准答案片段 $x_p$ 要被\textbf{屏蔽（mask）}，不参与梯度——因为那段不是模型自己写的，不该算它功劳。目标函数为：
\begin{equation*}
\mathrm{Objective}(\theta) = \mathbb{E}_{(x',\hat{y})\sim D_{\pi_\theta}}\big[\, r(x', \hat{y}) - \beta\,\log\big(\pi_\theta(\hat{y}\mid x') / \pi_{\mathrm{ref}}(\hat{y}\mid x')\big)\big]
\tag{4}
\end{equation*}
其中 $\pi_{\mathrm{ref}}$ 是 PPO 里的参考模型（由 SFT 模型初始化），$\beta$ 控制 KL 惩罚强度——这一项就是前面知识卡说的"拴住策略模型别跑太远"。
\item \textbf{起点逐步前移（课程推进）。} 随训练进行，起始点 $s^*$ 慢慢往标准答案的\textbf{开头}移动。每个样本设一个阈值 $\rho$：当模型在该样本上"累计正确比例"超过 $\rho$，就把起点往前挪一格（给的提示更少、要补的更多）。到训练后期 $s^* = 0$，模型完全只靠需求 $x$ 自己写整段代码——这时就退化成了标准 RL，但已经"学会怎么走到终点"了。
\item \textbf{采样次数与条件语句挂钩。} 条件语句越多 $\Rightarrow$ 独立路径越多 $\Rightarrow$ 逻辑越复杂，就需要更多次采样。作者把第 $i$ 个样本的课程数设为 $\lceil \sqrt{E_i}\, \rceil$（$E_i$ 是其条件语句数），课程步长设为 $\lceil \lceil E_i \rceil / \lceil \sqrt{E_i}\,\rceil \rceil$ 而非 1，以此加速训练。
\end{itemize}

\textbf{CCCS 一句话总结：} 从靠近目标的状态出发探索很容易；从远离目标的状态出发很难，但只要能复用"已经学会怎么抵达目标"的状态，就会变容易。这正是课程学习的精髓。

\paragraph{组件二：FGO（细粒度优化）。} 作者观察到：代码生成和 Atari 这类 RL 任务不同——\textbf{可以明确排除一批和奖励无关的动作}。既然编译器反馈只跟"被执行到的代码片段"有关，那就在计算损失时，把\textbf{没被单元测试执行到的 token 屏蔽掉}，只用执行到的 token 来更新策略模型。这样优势估计就只落在"真正影响结果"的代码上，优化更精准。

\paragraph{完整算法（Algorithm 1）。} CCCS 与 FGO 合在一起的训练循环大致是：对每个样本，先按当前课程算出起点位置 \texttt{pos} $= s_i \times c_i$，把"需求 + 标准答案前 \texttt{pos} 个 token"拼成新上下文 $x_i'$（CCCS）；采样轨迹、用式 (3) 算奖励；再算出"未执行片段的掩码矩阵" $m_{ij}$（执行到记 1、否则记 0，FGO）；用 PPO 更新模型；最后用滑动平均更新每个样本的通过率 $\rho_i$，一旦超过阈值 $\rho_t$ 就把课程 $c_i$ 往前推进一阶。如此循环。

% ---- 方法图解 ----
\begin{figure}[H]\centering
  \includegraphics[width=\linewidth]{cand_p04_k00.pdf}
  \caption{StepCoder 方法总览（原文 Figure 2 的矢量重制）。上半部分是普通的"编译器反馈强化学习（Vanilla RL）"；下半部分是 StepCoder。}
  \label{fig:overview}
\end{figure}

\paragraph{图~\ref{fig:overview} 在说什么（数据怎么流动）。}
\begin{itemize}[leftmargin=2em]
\item \textbf{上排（Vanilla RL）：} 人类需求 $\to$ LLM 生成整段代码 $\to$ 编译器跑测试给出一个 Reward。代码很长、只在最后拿一次稀疏奖励，所以探索很难。
\item \textbf{下排（StepCoder）：} 左侧大箭头是 \textbf{CCCS} 的课程进程。Step~1 时，输入 = 需求 + \emph{一大段}标准答案（图中实线文档表示"给定的提示"，虚线框表示"待补全的部分"），模型只做\textbf{代码补全（Code Completion）}，很容易成功；越往下走，给的提示越少、虚线区域越大，难度递增；到 \textbf{Final} 行，提示几乎为空，模型要直接做完整的\textbf{代码生成（Code Generation）}。
\item \textbf{右侧的 Mask 与 FGO：} 每行生成的代码送进编译器后，会得到一个掩码——图中\textcolor{green!60!black}{绿色}表示被执行到的代码、\textcolor{red}{红色}表示没被执行到的代码。\textbf{FGO 只对绿色（执行到的）token 计算梯度 $\sum_t \log \pi(y_t\mid y_{1:t-1}, x)$}，红色部分被屏蔽掉，从而实现细粒度优化。
\end{itemize}

% ======================================================================
\section{实验解读}
% ======================================================================

\subsection{数据集：APPS+ 怎么来的}

作者发现现有开源数据集里只有 APPS 勉强够格做 RL 训练，但它有不少脏数据：缺输入/输出/标准答案、标准答案编译不了或跑不起来、输出对不上等。于是他们做了一轮细致清洗：剔除缺字段的实例；统一输入输出格式以保证执行和比较时的确定性；对每个实例\textbf{跑单元测试 + 人工分析}，删掉代码不全/不相关、语法错误、API 误用、缺依赖库的；输出对不上的则人工复核题面、修正期望输出或删除。

最终得到 \textbf{APPS+}，共 \textbf{7{,}456} 个实例，按难度分三档：Introductory（入门，2{,}850）、Interview（面试，4{,}020）、Competition（竞赛，586）。平均每题描述 255.3 词、代码 21.9 行、附带 3 个单元测试，并标注了条件语句的起止位置。

\subsection{基准、基线与设置}

\begin{itemize}[leftmargin=2em]
\item \textbf{评测基准：} 自建的 APPS+，外加两个公认基准 MBPP（Mostly Basic Programming Problems，974 题入门级）和 HumanEval（164 道手写 Python 题）。MBPP/HumanEval 用\textbf{零样本（zero-shot）}评测——模型只在 APPS+ 上微调，然后直接去这两个基准上考。
\item \textbf{基线模型：} 基座模型（CodeLlama 系列、DeepSeek-Coder-Base）、SFT 模型（StarCoder、CodeLlama-Instruct、WizardCoder、DeepSeek-Coder-Instruct 等），以及 RL 类方法 Vanilla PPO、PPOCoder、RLTF。为公平比较，所有 RL 方法都用同一个基座 \textbf{DeepSeek-Coder-Instruct-6.7B}。
\item \textbf{评价指标：} Pass@k——对每个需求生成 $k$ 份代码，只要有一份通过全部单元测试就算"过"。本文主要看 \textbf{pass@1}。
\item \textbf{关键超参：} PPO 阶段策略模型学习率 $5\mathrm{e}{-7}$、critic 模型 $1.5\mathrm{e}{-6}$；每个样本采样 16 份；KL 惩罚系数 $\beta = 0.05$、clip 值 0.8；最大输出长度 1024。
\end{itemize}

\subsection{主要结果}

\paragraph{在 APPS+ 上（表~\ref{tab:apps}）。} RL 类方法整体优于基座模型和 SFT 模型，说明用编译器反馈做 RL 确实能更好地"导航"输出空间。\textbf{StepCoder 全面领先所有基线}，Overall 达 36.1，比最强的 RL 基线 RLTF（32.7）高出 3.4 个百分点；在竞赛级（Competition）这种最难的题上也拿到最高分 8.6。值得注意：在同一基座上做 SFT（"SFT on APPS+"，Overall 29.8）几乎没比基座（29.2）涨多少，反衬出\textbf{编译器反馈 RL 比单纯的下一 token 预测更能提升代码质量}。

\begin{table}[H]\centering\caption{在 APPS+ 上的 pass@1 结果（节选自原文 Table 1）。RL 类方法统一以 DeepSeek-Coder-Instruct-6.7B 为基座。}
\label{tab:apps}
\small
\begin{tabular}{lccccc}
\toprule
模型 & 规模 & Introductory & Interview & Competition & Overall \\
\midrule
\multicolumn{6}{l}{\textit{基座模型（Base Models）}}\\
CodeLlama & 13B & 18.7 & 11.0 & 0.0 & 13.0 \\
CodeLlama-Python & 13B & 29.0 & 12.3 & 2.9 & 17.9 \\
DeepSeek-Coder-Base & 6.7B & 13.0 & 10.3 & 5.0 & 10.9 \\
\midrule
\multicolumn{6}{l}{\textit{监督微调模型（SFT Models）}}\\
StarCoder & 15.6B & 6.3 & 4.1 & 0.7 & 4.7 \\
CodeLlama-Instruct & 13B & 33.3 & 11.0 & 1.4 & 18.7 \\
WizardCoder-Python-V1.0 & 13B & 39.7 & 15.1 & 4.3 & 23.6 \\
DeepSeek-Coder-Instruct & 6.7B & 49.4 & 18.7 & 3.6 & 29.2 \\
SFT on APPS+ & 6.7B & 50.1 & 19.0 & 6.4 & 29.8 \\
\midrule
\multicolumn{6}{l}{\textit{RL 类方法（基座 = DeepSeek-Coder-Instruct-6.7B）}}\\
Vanilla PPO & 6.7B & 53.7 & 20.1 & 5.0 & 31.7 \\
PPOCoder & 6.7B & 54.4 & 20.3 & 6.4 & 32.1 \\
RLTF & 6.7B & 55.1 & 20.8 & 6.4 & 32.7 \\
\textbf{StepCoder（本文）} & 6.7B & \textbf{59.7} & \textbf{23.5} & \textbf{8.6} & \textbf{36.1} \\
\quad w/o CCCS & 6.7B & 58.7 & 21.7 & 7.1 & 34.6 \\
\quad w/o FGO & 6.7B & 58.4 & 23.3 & 8.6 & 35.5 \\
\bottomrule
\end{tabular}
\end{table}

\paragraph{消融实验（Ablation）。} 看表~\ref{tab:apps} 最后两行：去掉 CCCS（w/o CCCS）后 Overall 从 36.1 掉到 34.6，且竞赛级从 8.6 大跌到 7.1——印证 \textbf{CCCS 主要帮在最难的复杂题上}（因为它简化了复杂需求的探索）；去掉 FGO（w/o FGO）后 Overall 掉到 35.5——说明 \textbf{FGO 通过把编译器反馈对准"被执行到的代码"提升了单测通过率}。两个组件都有正贡献。

\paragraph{在 MBPP 与 HumanEval 上（表~\ref{tab:mbpp}）。} 模型在 APPS+ 上训练后零样本迁移到这两个基准，\textbf{StepCoder 在两个基准上都最强}（HumanEval 78.7、MBPP 67.0）。作者还担心"是不是训练集和测试集重叠才涨的"，于是计算了 APPS+ 与两基准的代码行重叠率：只有 0.2\%（对 MBPP）和 7.1\%（对 HumanEval）的题目有超过一半的行能在对方里匹配上——重叠极小，说明提升是真的泛化能力，主要来自"改善了 RL 的探索问题"。另外一个有意思的发现：\textbf{只在单一数据集上做 SFT 反而会掉点}（"SFT on APPS+" 在两个基准上都明显下降），因为这会损害模型的指令跟随和泛化能力；而 RL 类方法则能在未见任务上提升表现。

\begin{table}[H]\centering\caption{在 MBPP 与 HumanEval 上的 pass@1 结果（零样本，原文 Table 2）。所有模型均在 APPS+ 上微调后评测。}
\label{tab:mbpp}
\small
\begin{tabular}{lcc}
\toprule
模型（6.7B） & HumanEval & MBPP \\
\midrule
DeepSeek-Coder-Instruct & 78.0 & 64.2 \\
SFT on APPS+ & 55.5 & 54.8 \\
Vanilla PPO & 78.0 & 65.0 \\
PPOCoder & 76.8 & 63.8 \\
RLTF & 76.8 & 65.2 \\
\textbf{StepCoder（本文）} & \textbf{78.7} & \textbf{67.0} \\
\bottomrule
\end{tabular}
\end{table}

\paragraph{按单测结果细分（原文 Figure 4）。} 作者进一步分析"没通过测试"的失败案例，把失败分成"编译错误（Reward $=-1$）"和"运行时错误与失败（Reward $=-0.6$ 或 $-0.3$）"两类。结论：StepCoder \textbf{显著降低了编译错误的比例}（在面试级、竞赛级题上尤其明显），说明它更容易写出"至少能跑"的代码；不过所有 LLM 都更容易栽在运行时错误上，StepCoder 也不例外（只是相对更低）。作者由此建议：\textbf{未来研究应重点降低运行时错误}，这能进一步提升代码质量与通过率。

\subsection{相关工作（Related Work）速览}

\begin{itemize}[leftmargin=2em]
\item \textbf{代码生成的 LLM：} 从预训练模型（CodeGPT、PanGu-Coder、SantaCoder、CodeGeex、Phi-1.5）到 SFT 模型（CodeX、StarCoder、WizardCoder、CodeLlama-Instruct、DeepSeek-Coder）；以及把 RL 引入代码生成的 CodeRL、PPOCoder、RLTF——但它们的奖励空间仍然稀疏，复杂任务探索依旧难。
\item \textbf{RL 中的探索：} 用 SFT 模型初始化策略来改善探索是常见做法（本文也用了），但对复杂需求还不够；过程监督奖励模型（process-supervised reward model）能给逐步奖励，但需要标注大量偏好数据；也有工作用"逐步变难的起始状态"构建课程（多见于游戏、机器人）。\textbf{本文的差异化}在于：结合软件工程特征，用\textbf{条件语句}来动态决定起始状态，并引入 FGO 利用\textbf{代码覆盖信息}做细粒度优化。
\end{itemize}

% ======================================================================
\section{核心 takeaways 与对 AI 应用开发的启示}
% ======================================================================

\begin{enumerate}[leftmargin=2em]
\item \textbf{"能不能跑通"是最硬的监督信号。} 对代码、SQL、配置、调用工具的参数这类\textbf{可自动验证}的任务，与其堆人工标注，不如把"执行结果对不对"做成自动奖励——这就是当下很火的 RLVR（基于可验证奖励的强化学习）思路。StepCoder 是把它在代码场景做扎实的一个范例。
\item \textbf{探索难，就上"课程"。} 当目标序列很长、奖励很稀疏时，让模型"从接近答案的地方起步、逐步退到从零开始"是个通用且便宜的技巧。做 agent / 长链路任务时，可以借鉴"先给一半轨迹当提示、再慢慢撤掉"的课程式训练。
\item \textbf{信用分配要对准"真正起作用的部分"。} FGO 的本质是\textbf{只对影响结果的 token 算梯度}。在你自己的 RL/微调里，如果能识别出"哪些输出真正决定了奖励"（如被执行到的代码、被实际调用的工具、被引用的检索片段），就该把优化聚焦在那里，别让无关 token 稀释信号。
\item \textbf{数据质量决定 RL 上限。} 作者花大力气清洗出 APPS+（人工核验单测、统一 I/O 格式、保证确定性比较）。一句话：\textbf{脏的单元测试 = 错的奖励 = 训歪的模型}。自动验证类项目里，先把"验证器/测试集"本身做对，比调模型更重要。
\item \textbf{单一数据集 SFT 可能伤泛化。} 实验里"只在 APPS+ 上做 SFT"反而让 HumanEval/MBPP 掉点，而 RL 方法能在未见任务上提升。启示：垂直数据上猛灌 SFT 要小心\textbf{灾难性遗忘 / 指令跟随退化}，RL（带 KL 约束）相对更稳。
\item \textbf{知道短板在哪。} StepCoder 主要降的是\textbf{编译错误}，对\textbf{运行时错误}帮助有限。做评测和迭代时，按错误类型拆开看，能更精准地知道下一步该补哪块。
\end{enumerate}

% ======================================================================
\section*{参考文献（节选最关键几条，原文不译）}
% ======================================================================
\begin{enumerate}[leftmargin=2em]
\item Hendrycks et al. 2021. Measuring coding challenge competence with APPS. NeurIPS Datasets and Benchmarks 2021. \textit{(APPS 数据集，本文 APPS+ 的来源)}
\item Austin et al. 2021. Program synthesis with large language models. arXiv:2108.07732. \textit{(MBPP 基准)}
\item Chen et al. 2021. Evaluating large language models trained on code. arXiv:2107.03374. \textit{(HumanEval 基准 \& Pass@k 指标)}
\item Le et al. 2022. CodeRL: Mastering code generation through pretrained models and deep reinforcement learning. NeurIPS 35. \textit{(用单测做奖励的 actor-critic 方法)}
\item Shojaee et al. 2023. Execution-based code generation using deep reinforcement learning (PPOCoder). arXiv:2301.13816.
\item Liu et al. 2023. RLTF: Reinforcement learning from unit test feedback. arXiv:2307.04349.
\item Schulman et al. 2017. Proximal policy optimization algorithms (PPO). arXiv:1707.06347.
\item Guo et al. 2024. DeepSeek-Coder: When the large language model meets programming. \textit{(本文 RL 实验的统一基座)}
\end{enumerate}

\vspace{1em}
\noindent\rule{\linewidth}{0.4pt}\\
\small\textit{本导读为转化性中文学习材料，仅用于学习理解；论文版权归原作者与 ACL 所有。原文与数据集/代码见上方 GitHub 链接。}

\end{document}
